\section{The matching algorithm}
\label{sec:matching}

The matching loop closes two data paths each Cauchy step: the worldtube
path (Cauchy state $\to$ characteristic boundary data,
Secs.~\ref{sec:n1}--\ref{sec:wtmap}) and the datum path (characteristic
$\psi_0$ $\to$ the physical boundary rows of Sec.~\ref{sec:bcrows},
Secs.~\ref{sec:retarded}--\ref{sec:cadence}). The full derivation chain
is summarized in Fig.~\ref{fig:dag}.

\subsection{Worldtube data map (node N1)}
\label{sec:n1}

\begin{equation}
g_{00} = -\alpha^2 + \gamma_{ij}\beta^i\beta^j,\qquad
g_{0i} = \gamma_{ij}\beta^j,\qquad
g_{ij} = \gamma_{ij};
\qquad
g^{00} = -\alpha^{-2},\quad
g^{0i} = \frac{\beta^i}{\alpha^{2}},\quad
g^{ij} = \gamma^{ij} - \frac{\beta^i\beta^j}{\alpha^{2}}.
\label{eq:fourmetric}
\end{equation}
ADM kinematics requires
$\partial_t\gamma_{ij} = -2\alpha K_{ij} + \beta^k\partial_k\gamma_{ij}
 + \gamma_{kj}\partial_i\beta^k + \gamma_{ik}\partial_j\beta^k$.
From $\gamma_{ij} = \chi^{-1}\tilde\gamma_{ij}$,
\begin{equation}
\partial_t\gamma_{ij}
 = \chi^{-1}\partial_t\tilde\gamma_{ij}
 - \chi^{-2}\tilde\gamma_{ij}\,\partial_t\chi ,
\label{eq:dtg-step1}
\end{equation}
and inserting \eqref{eq:dtchi}--\eqref{eq:dtgt} with
$K_{ij} = \chi^{-1}\big(\tilde A_{ij}+\tfrac13\tilde\gamma_{ij}K\big)$,
$\beta^k\partial_k\gamma_{ij} = \chi^{-1}\beta^k\partial_k\tilde\gamma_{ij}
 - \gamma_{ij}\,\beta^k\partial_k\ln\chi$,
\begin{equation}
\big[\partial_t\gamma_{ij}\big]_{\rm Z4c}
 - \big[\partial_t\gamma_{ij}\big]_{\rm ADM}
 = \gamma_{ij}\Big[\tfrac{2}{3}\,D_k\beta^k
 - \tfrac{2}{3}\,\partial_k\beta^k
 + \beta^k\partial_k\ln\chi\Big].
\label{eq:dtg-step2}
\end{equation}
\begin{nfblock}
\begin{equation}
\boxed{\;D_i\beta^i = \partial_k\beta^k
 - \tfrac{3}{2}\,\beta^k\,\partial_k\ln\chi\;}
\qquad\text{[V3: closure exact; naive }D=\partial_k\beta^k\text{ leaves
residual }+1\cdot\gamma_{ij}\beta^k\partial_k\ln\chi\text{]}.
\label{eq:divbeta}
\end{equation}
\end{nfblock}

\subsection{Live sampling and the $\ell=2$ data contract}
\label{sec:contract}

\begin{nfblock}
The Cauchy ADM state is interpolated to three geodesic spheres at
$r_{\rm wt}$ and $r_{\rm wt}\pm\delta r$ each step. With orthonormal
frame components $h_{\hat a\hat b} = (\gamma - \delta)_{\hat a\hat b}$
and $\partial_t h_{\hat a\hat b} = -2K_{\hat a\hat b}$ (linear order;
$\alpha = 1 + O(h)$, the correction entering beyond linear), the
$(\ell,m)=(2,0)$ contract scalars are the shape projections
\begin{equation}
h_{\rm TT} = \Big\langle \tfrac12(h_{\hat\theta\hat\theta}
 - h_{\hat\phi\hat\phi})\Big\rangle_{s^2},
\quad
h_{rr} = \big\langle h_{\hat r\hat r}\big\rangle_{P},
\quad
\varh = \big\langle h_{\hat\theta\hat\theta}
 + h_{\hat\phi\hat\phi}\big\rangle_{P},
\quad
h_{r\theta} = \big\langle r\,h_{\hat r\hat\theta}\big\rangle_{sc},
\label{eq:contract}
\end{equation}
with $\langle f\rangle_X = \int f X\,d\Omega / \int X^2\,d\Omega$ for
$X \in \{s^2, P, sc\}$, the time derivatives obtained from the same
projections of $-2K$, and $\partial_r\varh$ by the centered difference
across the sphere triplet. At $\ell=2$ the angular derivatives close
analytically on the same scalars,
\begin{equation}
\big\langle\partial_\theta \varh\big\rangle_{sc} = -6\,\varh,
\qquad
\big\langle\partial_\theta \partial_r\varh\big\rangle_{sc}
 = -6\,\partial_r\varh,
\qquad
\big\langle(\partial_\theta + \cot\theta)\,h_{r\theta}\big\rangle_{P}
 = h_{r\theta},
\label{eq:dthclosures}
\end{equation}
so the contract is seven numbers per step:
$(h_{\rm TT}, \partial_t h_{\rm TT}, h_{rr}, \varh, \partial_r\varh,
 h_{r\theta}, \partial_t h_{r\theta})$.
\end{nfblock}

\subsection{The anchored worldtube map (six rows $+\,\beta$)}
\label{sec:wtmap}

\begin{nfblock}
The Cauchy coordinates near the worldtube do not satisfy the Bondi gauge
conditions; the radial gauge generator is fixed by areal-radius
consistency only up to functions of $(u,\theta)$. Anchoring the
generator to vanish \emph{at the worldtube}, with the angular generator
corrected so that the $g_{r\theta}$ gauge condition holds consistently
($\xi_a = \xi_{\rm old} + \partial_\theta(\zeta_u^{\rm wt})/r -
\text{const}$), makes the map from the contract
\eqref{eq:contract}--\eqref{eq:dthclosures} to the characteristic
boundary scalars \emph{local}:
\begin{align}
J_t^{\rm wt} &= h_{\rm TT},
&
H_t^{\rm wt} &= \partial_t h_{\rm TT},
&
U_t^{\rm wt} &= \frac{6\,\varh}{4\,r_{\rm wt}},
\nonumber\\
Q_t^{\rm wt} &= \frac{6\,\varh}{4} + \frac{6\,r_{\rm wt}\,
\partial_r\varh}{4}
 - \frac{3\,h_{r\theta}}{r_{\rm wt}} + \partial_t h_{r\theta},
&
W_t^{\rm wt} &= \frac{3\,h_{rr}}{4\,r_{\rm wt}}
 - \frac{\partial_r\varh}{4},
&
\beta_t^{\rm wt} &= \frac{h_{rr}}{4} - \frac{\varh}{8}
 - \frac{r_{\rm wt}\,\partial_r\varh}{8}
 + \frac{h_{r\theta}}{4\,r_{\rm wt}},
\label{eq:sixrowmap}
\end{align}
verified as an exact identity against the anchored-gauge solution
series (BigFloat-192 residual $<10^{-35}$; derivation artifact
\texttt{derive\_wtmap\_v2.jl}). The anchored gauge necessarily carries
$\beta_t \neq 0$ and non-decaying tails $a(u) + b(u)/r$ in $J$, $U$,
$W$; these are exactly the structures supported by the
$\beta$-corrected, unringed system \eqref{eq:unringed}, and they are
$\psi_0$-silent by \eqref{eq:psi0formula}. Independent anchoring of the
radial and angular generators violates the $g_{r\theta}$ condition and
is excluded [pinned: the uncorrected map fails the $\beta$ channel with
a characteristic $(4-n)$-ratio pattern].
\end{nfblock}

\subsection{Retarded matching geometry}
\label{sec:retarded}

\begin{nfblock}
The characteristic time is the retarded cone label, $u = t - r$ on the
flat background: the cone $u$ meets the worldtube worldline at Cauchy
time $t = u + r_{\rm wt}$ and the Cauchy outer boundary $r_B$ at
$t = u + r_B$. Hence
\begin{equation}
\text{tube sample at } t \;\longrightarrow\; \text{cone } u = t -
r_{\rm wt},
\qquad
\text{boundary datum at } t \;\longleftarrow\;
\psi_0^t\big(r_B\big)\Big|_{\,u = t - r_B},
\label{eq:conelabels}
\end{equation}
a causality lag of $r_B - r_{\rm wt}$ in $u$ that is always available
in lockstep evolution. The solver records the interior probe
\eqref{eq:proberules} at $r_B$ each substep and serves the lagged value
by linear interpolation in the recorded history, with the analytic
initial-data value as the pre-cone fallback (the run starts on the
quiescent cone $u_0 = -r_{\rm wt}$). Boundary data are held fixed
across the substeps between Cauchy samples --- an $O(\Delta t)$
coupling approximation. Measured against the published reference data
of Ref.~\cite{Ma:2023qjn} (Sec.~\ref{sec:tests}), this approximation,
together with the tube-inside-datum geometry, produces $O(30\%)$
deviations of the live boundary $\psi_0$ from the clean characteristic
solution in the pulse-arrival window; linear-in-time boundary-data
interpolation is the corresponding open refinement [O-M5-1, class (c)].
\end{nfblock}

\subsection{Injection path, cadence, and initialization}
\label{sec:cadence}

\begin{nfblock}
The served datum enters the Cauchy evolution through the Type-II boost
and the physical rows of Sec.~\ref{sec:bcrows}:
$\psi_0' = A^2\,\psi_0$ with
$A = (\alpha - \gamma_{ij}\beta^is^j)\,e^{-2\hat\beta}$
\eqref{eq:boost}, injected as \eqref{eq:teuk-inject}. Per Cauchy RK
stage at the boundary: (1) sample the contract \eqref{eq:contract}
(collective over ranks; off-rank sphere points are excluded by
ownership, the shape norms being global); (2) map by
\eqref{eq:sixrowmap}; (3) advance the characteristic state to
$u = t - r_{\rm wt}$ (root rank); (4) serve the lagged probe
\eqref{eq:conelabels} and broadcast; (5) apply the boundary rows. The
characteristic initial slice is the analytic datum (quiescent for the
matched tests; \eqref{eq:pulseid} for the pulse-injection test);
initialization order and the $\hat\beta$ circularity at $t=0$ remain
documented obligations [O-N6-3]. Setting the served datum to zero
reduces the entire loop to the absorbing scheme of
Ref.~\cite{Ruiz:2010qj} exactly [V6].
\end{nfblock}

\subsection{Derivation chain}
\label{sec:dag}

\begin{figure}[htbp]
\centering
\begin{tikzpicture}[
  node distance = 3.5mm and 4.5mm,
  box/.style  = {draw, rounded corners=1pt, align=center,
                 font=\scriptsize, inner sep=2.2pt},
  newbox/.style = {box, draw=blue!60!black, text=blue!60!black, thick},
  arr/.style  = {-{Latex[length=1.6mm]}, line width=0.5pt, gray!70},
]
% layer 1: foundations
\node[box] (z4c)  {Z4c system\\ \eqref{eq:dtchi}--\eqref{eq:constraints}};
\node[box, right=of z4c] (frame) {boundary frame\\ \eqref{eq:adm}--\eqref{eq:nullvecs}};
\node[box, right=of frame] (bs) {Bondi--Sachs\\ \eqref{eq:bsmetric}--\eqref{eq:ycomp}};
\node[box, right=of bs] (teuk) {exact Teukolsky\\ \eqref{eq:teuk-F}--\eqref{eq:teuk-psi0}, N12};
% layer 2
\node[newbox, below=of z4c] (modes) {ten modes\\ Tab.~\ref{tab:modes}, N4--N6};
\node[newbox, below=of frame] (boost) {boost $A$\\ \eqref{eq:boost}, N2};
\node[newbox, right=of boost] (dict) {dictionary\\ \eqref{eq:exact-dictionary}, N3};
\node[newbox, below=of bs, xshift=18mm] (scal) {$\ell{=}2$ scalars\\ \eqref{eq:tscalars}--\eqref{eq:modalfactors}};
\node[newbox, right=of scal] (teukbondi) {Bondi forms\\ \eqref{eq:teukbondi}--\eqref{eq:teukpsi0}};
% layer 3
\node[newbox, below=of modes] (rows) {CPBC/gauge rows + N13\\ \eqref{eq:bc-cpbc}--\eqref{eq:bc-gauge}};
\node[newbox, below=of boost] (phys) {physical rows\\ \eqref{eq:bc-phys}, \eqref{eq:teuk-inject}};
\node[newbox, below=of scal] (plain) {plain hierarchy\\ \eqref{eq:plainhierarchy}};
% layer 4
\node[newbox, below=of phys, xshift=-12mm] (contract) {$\ell{=}2$ contract\\ \eqref{eq:contract}--\eqref{eq:dthclosures}};
\node[newbox, right=of contract] (map) {anchored map\\ \eqref{eq:sixrowmap}};
\node[newbox, right=of map] (beta) {$\beta$-hierarchy\\ \eqref{eq:beta-hierarchy}, \eqref{eq:unringed}};
\node[newbox, right=of beta] (psi0) {$\psi_0$ probe\\ \eqref{eq:psi0formula}--\eqref{eq:proberules}};
% layer 5
\node[newbox, below=of map, xshift=10mm] (retard) {retarded matching\\ \eqref{eq:conelabels}};
% layer 6
\node[box, below=of retard, xshift=-16mm] (tests) {batteries: modes 3--6\\ Sec.~\ref{sec:tests}};
\node[box, right=of tests] (oracle) {SpECTRE oracle\\ Sec.~\ref{sec:tests}};
\node[box, right=of oracle] (repro) {paper-data overlays\\ Sec.~\ref{sec:tests}};
% edges
\draw[arr] (z4c) -- (modes);
\draw[arr] (frame) -- (modes);
\draw[arr] (frame) -- (boost);
\draw[arr] (frame) -- (dict);
\draw[arr] (bs) -- (scal);
\draw[arr] (teuk) -- (teukbondi);
\draw[arr] (scal) -- (teukbondi);
\draw[arr] (modes) -- (rows);
\draw[arr] (modes) -- (phys);
\draw[arr] (boost) -- (phys);
\draw[arr] (dict) -- (phys);
\draw[arr] (scal) -- (plain);
\draw[arr] (z4c) to[out=200,in=170] (contract);
\draw[arr] (contract) -- (map);
\draw[arr] (plain) -- (beta);
\draw[arr] (map) -- (beta);
\draw[arr] (beta) -- (psi0);
\draw[arr] (teukbondi) -- (psi0);
\draw[arr] (map) -- (retard);
\draw[arr] (psi0) -- (retard);
\draw[arr] (phys) to[out=270,in=160] (retard);
\draw[arr] (retard) -- (tests);
\draw[arr] (teukbondi) to[out=300,in=30] (oracle);
\draw[arr] (retard) -- (oracle);
\draw[arr] (retard) -- (repro);
\draw[arr] (rows) to[out=270,in=180] (tests);
\end{tikzpicture}
\caption{Derivation chain of the Z4c-CCM formulation. Black nodes:
foundations (literature content and the exact reference chain); blue
nodes: machine-verified new content (ledger nodes named in the section
titles); bottom row: the verification instruments of
Sec.~\ref{sec:tests}.}
\label{fig:dag}
\end{figure}
